% Chapter 21: High-Precision Computation
% ch21_high_precision.tex

\chapter{High-Precision Computation: Pi, $e$, and $\gamma$}

\begin{historical}
Euler computed mathematical constants to extraordinary precision using only pen and paper.
His evaluation of $\pi$ to 150 digits and $e$ to many decimal places set standards that
would not be surpassed for centuries. In this chapter, we replicate Euler's spirit using
pure Python arbitrary-precision arithmetic --- no external libraries, no floating-point
rounding, just integer math at whatever precision we desire.

Euler's original methods included infinite series, continued fractions, and product formulas.
We implement several of these, plus the modern AGM (Arithmetic-Geometric Mean) algorithm
for $\pi$ and accelerated series for $\gamma$. The AGM method, discovered by Gauss but
implemented efficiently using Euler's ideas, converges quadratically and is used in
all modern high-precision $\pi$ computations.
\end{historical}

\section{Prerequisites}

To work through this chapter, you need:
\begin{itemize}
    \item \textbf{Mathematics:} Infinite series, continued fractions, basic understanding of convergence rates.
    \item \textbf{Programming:} Python integer arithmetic (arbitrarily large), iterative algorithms.
\end{itemize}

\section{Arbitrary-Precision Arithmetic}

Standard floating-point numbers in Python (and all programming languages) use IEEE 754
double-precision format, which gives about 15-17 significant decimal digits. For computing
constants to hundreds or thousands of digits, we need arbitrary-precision arithmetic.

Our approach stores numbers as scaled integers: a value $v$ is represented as
$V / 10^p$ where $V$ is an exact integer and $p$ is the precision. All arithmetic
operations are performed on integers, guaranteeing exact results (except for division,
which truncates to the chosen precision).

\subsection{The ScaledInt Class}

The \texttt{ScaledInt} class provides arbitrary-precision decimal arithmetic:

\begin{lstlisting}[caption={Basic high-precision arithmetic}, label=lst:hp_basic]
from euler import ScaledInt

# Create high-precision numbers
a = ScaledInt("3.141592653589793238462643383279", 30)
b = ScaledInt("2.718281828459045235360287471352", 30)

print(f"a = {a}")
print(f"b = {b}")
print(f"a + b = {a + b}")
print(f"a * b = {a * b}")
print(f"a / b = {a / b}")
print(f"sqrt(a) = {ScaledInt(2, 30).sqrt()}")
\end{lstlisting}

\subsection{Implementation Details}

Arithmetic operations on \texttt{ScaledInt} work as follows:
\begin{itemize}
    \item \textbf{Addition/Subtraction:} Convert to common precision, then add/subtract integers.
    \item \textbf{Multiplication:} Multiply integers, then divide by scale factor.
    \item \textbf{Division:} Multiply numerator by scale, then divide.
    \item \textbf{Square root:} Use Newton's method with integer arithmetic.
\end{itemize}

\subsection{Newton's Method for Square Roots}

To compute $\sqrt{N}$ to $p$ decimal places:
\begin{enumerate}
\item Start with initial guess $x_0 = \lfloor \sqrt{N} \rfloor \cdot 10^{p/2}$
\item Iterate: $x_{k+1} = \frac{1}{2}\left(x_k + \frac{N \cdot 10^p}{x_k}\right)$
\item Converges quadratically: digits of precision double each iteration
\end{enumerate}

\begin{lstlisting}[caption={High-precision square root}, label=lst:hp_sqrt]
from euler import ScaledInt

# Compute sqrt(2) to 50 digits
sqrt2 = ScaledInt(2, 50).sqrt()
print(f"sqrt(2) = {sqrt2}")

# Verify by squaring
print(f"sqrt(2)^2 = {sqrt2 * sqrt2}")
print(f"Difference from 2: {2 - sqrt2 * sqrt2}")
\end{lstlisting}

\section{Computing $\pi$}

\subsection{Machin's Formula}

Machin's formula (1706), contemporaneous with Euler's early work:
\[
\frac{\pi}{4} = 4 \arctan\frac{1}{5} - \arctan\frac{1}{239}
\]

Using the Taylor series $\arctan(x) = x - x^3/3 + x^5/5 - \cdots$, the terms
$1/5$ and $1/239$ converge rapidly. Machin used this to compute 100 digits of $\pi$.

\subsection{Convergence Analysis}

The Taylor series for $\arctan(x)$ converges for $|x| \leq 1$. The error after $n$ terms is approximately:
\[
R_n \approx \frac{x^{2n+1}}{2n+1}
\]

For $x = 1/5$, each term reduces the error by a factor of $1/25$, so we gain approximately 1.4 decimal digits per term.
For $x = 1/239$, convergence is even faster (about 4.4 decimal digits per term).

Machin's formula therefore gains about $4 \times 1.4 = 5.6$ digits per iteration from the first term,
minus about 4.4 digits from the second term, for a net gain of approximately 1.2 digits per term.

\begin{lstlisting}[caption={Pi via Machin's formula}, label=lst:pi_machin]
from euler import pi_machin

print(f"pi (50 digits): {pi_machin(50)}")
print(f"pi (100 digits): {pi_machin(100)}")
print(f"pi (200 digits): {pi_machin(200)}")
\end{lstlisting}

\subsection{Other Machin-like Formulas}

Euler discovered many related formulas. Some notable examples:

\begin{align*}
\frac{\pi}{4} &= 5 \arctan\frac{1}{7} + 2 \arctan\frac{3}{79} \quad \text{(Euler, 1748)} \\
\frac{\pi}{4} &= \arctan\frac{1}{2} + \arctan\frac{1}{3} \quad \text{(valid but slow)} \\
\frac{\pi}{4} &= 2 \arctan\frac{1}{3} + \arctan\frac{1}{7} \quad \text{(Hutton, 1789)}
\end{align*}

\begin{lstlisting}[caption={Euler's pi formula}, label=lst:pi_euler]
from euler import pi_machin

# Euler's formula: pi/4 = 5*arctan(1/7) + 2*arctan(3/79)
print(f"pi (100 digits) via Euler: {pi_machin(100, formula='euler')}")
\end{lstlisting}

\subsection{The AGM Algorithm}

The Brent--Salamin (Gauss--Legendre) algorithm converges quadratically:
each iteration doubles the number of correct digits. This is dramatically faster
than Machin's formula.

\subsubsection{Algorithm Description}

Starting with $a_0 = 1$, $b_0 = 1/\sqrt{2}$, $t_0 = 1/4$, $p_0 = 1$:
\begin{align*}
a_{n+1} &= \frac{a_n + b_n}{2} \\
b_{n+1} &= \sqrt{a_n b_n} \\
t_{n+1} &= t_n - p_n(a_n - a_{n+1})^2 \\
p_{n+1} &= 2p_n
\end{align*}

Then $\pi \approx \frac{(a_n + b_n)^2}{4t_n}$.

\subsubsection{Convergence Rate}

The AGM converges quadratically: if $a_n$ and $b_n$ agree to $d$ decimal digits,
then $a_{n+1}$ and $b_{n+1}$ agree to approximately $2d$ digits. This means that
after $k$ iterations, we have approximately $2^k$ correct digits.

For 50 digits, we need only about $\log_2(50) \approx 6$ iterations!

\begin{lstlisting}[caption={Pi via AGM}, label=lst:pi_agm]
from euler import pi_gauss_agm

print(f"pi (AGM, 50 digits): {pi_gauss_agm(50)}")
print(f"pi (AGM, 100 digits): {pi_gauss_agm(100)}")
\end{lstlisting}

After just 5 iterations, the AGM gives 50+ correct digits of $\pi$.

\subsection{AGM Convergence Verification}

\begin{lstlisting}[caption={Verifying AGM convergence}, label=lst:agm_verify]
from euler import ScaledInt

def agm_iterations(a0, b0, n):
    """Compute AGM iterations and show convergence."""
    a, b = a0, b0
    t = ScaledInt(1, n) / 4
    p = ScaledInt(1, n)
    
    print(f"AGM convergence to {n} digits:")
    print(f"{'Iter':>4} | {'a':^30} | {'b':^30} | {'diff':^15}")
    print("-" * 80)
    
    for i in range(min(10, n//10 + 2)):
        a_new = (a + b) / 2
        b_new = (a * b).sqrt()
        t = t - p * (a - a_new) ** 2
        p = 2 * p
        a, b = a_new, b_new
        
        # Compute pi estimate
        pi_est = (a + b) ** 2 / (4 * t)
        diff = abs(a - b)
        print(f"{i:>4} | {str(a)[:30]:^30} | {str(b)[:30]:^30} | {str(diff)[:15]:^15}")
        
        if diff < ScaledInt(1, n) * ScaledInt(10, -n):
            break
    
    pi_final = (a + b) ** 2 / (4 * t)
    print(f"\nFinal pi estimate: {pi_final}")
    return pi_final

agm_iterations(ScaledInt(1, 100), ScaledInt(1, 100).sqrt(2), 100)
\end{lstlisting}

\section{Computing $e$}

\subsection{Series Method}

Euler's original series:
\[
e = \sum_{n=0}^{\infty} \frac{1}{n!} = 1 + 1 + \frac{1}{2} + \frac{1}{6} + \cdots
\]

Each term adds approximately one decimal digit of precision (since $n!$ grows faster than $10^n$).
The error after $N$ terms is bounded by $\frac{1}{N \cdot N!}$.

\subsection{Continued Fraction}

Euler also discovered the continued fraction for $e$:
\[
e = [2; 1, 2, 1, 1, 4, 1, 1, 6, 1, 1, 8, \ldots]
\]
This converges much faster than the series. The $n$-th convergent has error approximately $1/n!$.

\begin{lstlisting}[caption={Computing e}, label=lst:compute_e]
from euler import e_series, e_continued_fraction

print(f"e (series, 50 digits): {e_series(50)}")
print(f"e (continued fraction, 50 digits): {e_continued_fraction(50)}")
\end{lstlisting}

\subsection{Comparing Methods}

\begin{lstlisting}[caption={Comparing e computation methods}, label=lst:e_compare]
from euler import e_series, e_continued_fraction
import time

print("Comparing e computation methods:")
for digits in [20, 50, 100, 200]:
    start = time.time()
    e1 = e_series(digits)
    time1 = time.time() - start
    
    start = time.time()
    e2 = e_continued_fraction(digits)
    time2 = time.time() - start
    
    print(f"  {digits} digits: series={time1:.4f}s, CF={time2:.4f}s")
    print(f"    Match: {e1 == e2}")
\end{lstlisting}

\section{Computing Euler's Constant $\gamma$}

The Euler--Mascheroni constant is the hardest of the three to compute at high precision,
because no simple closed-form series exists. We use the identity:
\[
\gamma = \sum_{k=1}^{\infty} \left(\frac{1}{k} - \ln\left(1 + \frac{1}{k}\right)\right)
\]

combined with Euler--Maclaurin acceleration. The term $\ln(1+1/k)$ is computed using
the identity $\ln(1+1/k) = 2\,\mathrm{atanh}(1/(2k+1))$, which converges quadratically.

\subsection{Euler--Maclaurin Acceleration}

The asymptotic expansion gives:
\[
\gamma = \sum_{k=1}^{N} \left(\frac{1}{k} - \ln\left(1 + \frac{1}{k}\right)\right) + \frac{1}{2N} - \frac{1}{12N^2} + \frac{1}{120N^4} - \cdots
\]

By choosing $N$ appropriately and including correction terms, we can compute $\gamma$ to high precision.

\begin{lstlisting}[caption={Computing gamma}, label=lst:compute_gamma]
from euler import euler_gamma

print(f"gamma (20 digits): {euler_gamma(20)}")
print(f"gamma (30 digits): {euler_gamma(30)}")
print(f"gamma (50 digits): {euler_gamma(50)}")
\end{lstlisting}

\subsection{Alternative Series for Gamma}

Several other series converge faster than the basic definition:

\begin{align*}
\gamma &= \sum_{n=2}^{\infty} (-1)^n \frac{\zeta(n)}{n} \\
\gamma &= \int_0^1 \left(\frac{1}{\ln x} + \frac{1}{1-x}\right) dx \\
\gamma &= \lim_{n \to \infty} \left(\sum_{k=1}^n \frac{1}{k} - \ln n\right)
\end{align*}

\begin{lstlisting}[caption={Alternative gamma computation}, label=lst:gamma_alt]
from euler import ScaledInt
import math

def gamma_zeta_series(n_terms, precision):
    """Compute gamma using zeta series."""
    gamma = ScaledInt(0, precision)
    for k in range(2, n_terms + 2):
        zeta_k = zeta_high_precision(k, precision)
        term = ((-1) ** k) * zeta_k / k
        gamma = gamma + term
    return gamma

def gamma_integral_approx(n_steps, precision):
    """Approximate gamma integral."""
    # Using trapezoidal rule on [0,1]
    h = ScaledInt(1, precision) / n_steps
    total = ScaledInt(0, precision)
    
    for i in range(n_steps):
        x = h * (i + ScaledInt(0.5, precision))
        f = (ScaledInt(1, precision) / x.ln() + 
             ScaledInt(1, precision) / (ScaledInt(1, precision) - x))
        total = total + f * h
    
    return total
\end{lstlisting}

\section{Zeta Values}

High-precision $\pi$ enables high-precision zeta values at even integers:

\begin{lstlisting}[caption={High-precision zeta values}, label=lst:zeta_hp]
from euler import zeta_high_precision

print(f"zeta(2) = {zeta_high_precision(2, 50)}")
print(f"zeta(4) = {zeta_high_precision(4, 50)}")
print(f"zeta(6) = {zeta_high_precision(6, 50)}")
\end{lstlisting}

For odd integers, the computation is more challenging. The series $\sum n^{-s}$ converges
too slowly for direct computation. Accelerated series (such as Euler's transformation)
are required.

\subsection{Euler's Acceleration Technique}

Euler developed a powerful series acceleration method. For a series $\sum a_n$,
define the Euler transform:
\[
\sum_{n=0}^{\infty} a_n = \sum_{n=0}^{\infty} \frac{(-1)^n}{2^{n+1}} \sum_{k=0}^{n} \binom{n}{k} a_k
\]

\begin{lstlisting}[caption={Euler acceleration for zeta(3)}, label=lst:euler_accel_zeta3]
from euler import ScaledInt

def euler_accelerated_zeta3(precision, terms=100):
    """Compute Ap\'{e}ry's constant using Euler acceleration."""
    # Direct series is too slow, use acceleration
    a = [ScaledInt(1, precision) / (k + 1)**3 for k in range(terms)]
    
    result = ScaledInt(0, precision)
    for n in range(terms):
        inner = ScaledInt(0, precision)
        for k in range(n + 1):
            binom = math.comb(n, k)
            inner = inner + ScaledInt(binom, precision) * a[k]
        sign = (-1) ** n
        weight = sign / (2 ** (n + 1))
        result = result + weight * inner
    
    return result

# Note: This is a simplified illustration
# Actual implementation uses more sophisticated acceleration
\end{lstlisting}

\section{Computational Complexity Analysis}

\subsection{Machin's Formula}

Computing $\pi$ to $d$ decimal digits using Machin's formula requires $O(d^2)$ integer operations,
since each arctangent term requires $O(d)$ iterations and each iteration involves $O(d)$-bit arithmetic.

\subsection{AGM Algorithm}

The AGM algorithm requires $O(\log d)$ iterations, with each iteration involving $O(M(d))$ operations
where $M(d)$ is the cost of multiplying $d$-bit numbers. Using fast multiplication, the total
complexity is $O(M(d) \log d)$.

\subsection{Euler's Constant}

Computing $\gamma$ to $d$ digits using the accelerated series requires $O(d)$ terms,
each involving $O(d)$-bit arithmetic. The total complexity is $O(d^2)$.

\section{Fast Multiplication Algorithms}

For very high precision computations, the choice of multiplication algorithm matters.

\subsection{Schoolbook Multiplication}

Standard $O(d^2)$ algorithm suitable for moderate precision:

\begin{lstlisting}[caption={Schoolbook multiplication}, label=lst:schoolbook_mul]
def schoolbook_multiply(a, b, precision):
    """Multiply two high-precision numbers stored as scaled integers."""
    scale = 10 ** precision
    result = (a * b) // scale
    return result
\end{lstlisting}

\subsection{Karatsuba Multiplication}

Karatsuba's algorithm achieves $O(d^{\log_2 3}) \approx O(d^{1.585})$:

\begin{lstlisting}[caption={Karatsuba multiplication}, label=lst:karatsuba_mul]
def karatsuba_multiply(x, y, precision):
    """Karatsuba fast multiplication for scaled integers."""
    if x < 10**6 or y < 10**6:
        return (x * y) // (10 ** precision)
    
    n = max(len(str(x)), len(str(y))) // 2
    m = 10 ** n
    
    a, b = x // m, x % m
    c, d = y // m, y % m
    
    ac = karatsuba_multiply(a, c, precision)
    bd = karatsuba_multiply(b, d, precision)
    ab_cd = karatsuba_multiply(a + b, c + d, precision) - ac - bd
    
    return (ac * m * m + ab_cd * m + bd) // (10 ** precision)
\end{lstlisting}

\subsection{FFT-based Multiplication}

For extreme precision, FFT-based multiplication achieves near-linear complexity:

\begin{lstlisting}[caption={FFT multiplication concept}, label=lst:fft_mul]
def fft_multiply_concept(a, b, precision):
    """
    Conceptual FFT-based multiplication.
    In practice, use libraries like GMP for production code.
    """
    # Convert to base-2^30 representation for FFT
    # Apply FFT, pointwise multiply, inverse FFT
    # Reconstruct result
    pass  # Implementation omitted for brevity
\end{lstlisting}

\section{Modern Relevance}

High-precision computation is essential in: verifying the Riemann Hypothesis (zeta zeros
computed to trillions of digits), testing mathematical constants for normality, cryptographic
key generation, and benchmarking computer hardware. The Chudnovsky algorithm (used to compute
$\pi$ to trillions of digits) is a descendant of Euler's series methods.

\subsection{Pi Record Computations}

Some notable $\pi$ record computations:
\begin{itemize}
\item 1949: ENIAC computer computed 2037 digits
\item 1989: Chudnovsky brothers computed 1 billion digits
\item 2021: Google computed 100 trillion digits using cloud computing
\end{itemize}

\section{Frequently Asked Questions}

\begin{itemize}
\item \textbf{Q: Why not just use Python's built-in float?} \\
Python floats are IEEE 754 doubles with ~15-17 significant digits. For computing $\pi$ to 100 digits or verifying that $\zeta(3)$ is irrational, we need exact arithmetic. Our \texttt{ScaledInt} type stores every digit exactly.

\item \textbf{Q: How does the AGM algorithm work?} \\
Starting with $a_0=1$, $b_0=1/\sqrt{2}$, $t_0=1/4$, $p_0=1$, the iterations
$a_{n+1}=(a_n+b_n)/2$, $b_{n+1}=\sqrt{a_nb_n}$, $t_{n+1}=t_n-p_n(a_n-a_{n+1})^2$,
$p_{n+1}=2p_n$ converge quadratically to values from which $\pi$ is extracted.

\item \textbf{Q: Why is $\gamma$ harder to compute than $\pi$ and $e$?} \\
Unlike $\pi$ and $e$, $\gamma$ has no known closed-form expression in terms of other constants, no simple continued fraction, and no rapid-converging series. The best known methods converge as $O(1/N)$ even with acceleration.

\item \textbf{Q: What is the fastest known algorithm for $\pi$?} \\
The Chudnovsky algorithm (1988), which converges 14 digits per term, is the fastest. It is based on Ramanujan's series and is used by virtually all $\pi$-record computations.

\item \textbf{Q: Can this compute $\pi$ to a million digits?} \\
Yes, in principle. The AGM algorithm would need about $\log_2(10^6) \approx 20$ iterations. However, the Python implementation is not optimized for extreme precision --- production code uses specialized libraries like GMP.

\item \textbf{Q: How does the continued fraction for $e$ work?} \\
The continued fraction $e = [2; 1, 2, 1, 1, 4, 1, 1, 6, 1, 1, 8, \ldots]$ has a regular pattern. Each convergent $p_n/q_n$ gives an approximation with error $|e - p_n/q_n| < 1/q_n^2$.

\item \textbf{Q: What is the relationship between Machin's formula and Euler's work?} \\
Euler studied Machin's formula extensively and discovered many related formulas. He showed that if $\arctan(1/a) + \arctan(1/b) = \arctan(1/c)$, then $c = (ab-1)/(a+b)$. This allows generating infinitely many such formulas.

\item \textbf{Q: What is the Chudnovsky algorithm?} \\
The Chudnovsky algorithm uses the hypergeometric series:
\[
\frac{1}{\pi} = 12 \sum_{k=0}^{\infty} \frac{(-1)^k (6k)! (13591409 + 545140134k)}{(3k)!(k!)^3 640320^{3k+3/2}}
\]
It converges at about 14 digits per term, making it extremely efficient for record computations.
\end{itemize}

\section{Exercises}

\begin{exercise}
Implement a high-precision computation of $\sqrt{2}$ using Newton's method. Compare the convergence rate with the AGM algorithm for $\pi$.
\end{exercise}

\begin{exercise}
Use the continued fraction for $e$ to compute $e$ to 100 digits. How many terms are needed?
\end{exercise}

\begin{exercise}
Verify that $\zeta(2) = \pi^2/6$ to 30 digits using your high-precision $\pi$ and zeta implementations.
\end{exercise}

\begin{exercise}
Derive Machin's formula from the identity $\arctan(1/2) + \arctan(1/3) = \arctan(1)$ and compute $\pi$ to 20 digits using this alternative formula.
\end{exercise}

\begin{exercise}
Show that the error in the AGM algorithm decreases quadratically by computing $\pi$ to increasing precision and tracking the number of correct digits after each iteration.
\end{exercise}

\begin{exercise}
Implement Euler's acceleration technique for the series $\sum 1/n^2$ and compare its convergence with the direct partial sum.
\end{exercise}

\begin{exercise}
Use the identity $\gamma = \int_0^1 \frac{1-x}{(1-x)\ln x} dx$ to compute $\gamma$ to 10 decimal places using numerical integration.
\end{exercise}

\begin{exercise}
Prove that the continued fraction convergents of $e$ satisfy $|e - p_n/q_n| < 1/q_n^2$. How many convergents are needed to get 50 correct digits?
\end{exercise}

\begin{exercise}
Compare the computational complexity of computing $\pi$ to $d$ digits using Machin's formula versus the AGM algorithm. How many iterations does each require for $d = 100$?
\end{exercise}

\begin{exercise}
Prove that the AGM algorithm satisfies $\pi \approx \frac{(a_n + b_n)^2}{4t_n}$ by showing that $a_n \to \frac{\pi}{2M(1,1/\sqrt{2})}$ where $M$ is the arithmetic-geometric mean.
\end{exercise}

\begin{exercise}
Compute $\pi$ to 100 digits using both Machin's formula and the AGM algorithm. Compare the number of operations required for each method.
\end{exercise}

\begin{exercise}[Computational]
Implement the Chudnovsky algorithm and compute $\pi$ to 50 digits. Compare the number of terms required with Machin's formula.
\end{exercise}

\begin{exercise}
Show that the error in the series $\sum_{n=1}^{N} 1/n^2$ for $\zeta(2)$ is approximately $1/N$. Use Euler's acceleration to improve convergence.
\end{exercise}

\begin{exercise}
Implement Karatsuba multiplication and use it to speed up high-precision arithmetic. Compare performance with schoolbook multiplication for $d = 1000$ digits.
\end{exercise}

\begin{exercise}
Prove that if $\arctan(1/a) + \arctan(1/b) = \arctan(1/c)$, then $c = (ab-1)/(a+b)$. Use this to generate new Machin-like formulas for $\pi$.
\end{exercise}
